\color{unidarkgreen}{\subsection{Контроль отсутствия напряжения}}
\color{black}

\begin{enumerate}

\item
Контроль отсутствия напряжения своей стороны (КОН СШ) и низкой стороны (КОН НН) используется в качестве дополнительного разрешающего критерия в логике ускорения выключателя. 

\item
Уставки КОН представлены в таблице \ref{app_set:kon}. На рисунке \ref{pic_kon:KON} приведен алгоритм работы КОН. 

\begin{figure}[H]
\centering
\includegraphics[width=0.8\textwidth,height=0.8\textheight,keepaspectratio]{\fbpath/07.02 РЗАТ/051.1751 КОН/img/KON.pdf}
\captionof{figure}{Функциональная схема КОН}\label{pic_kon:KON}
\end{figure}

\item
Орган, контролирующий отсутствие напряжения низкой стороны (КОН НН), контролирует минимальное линейное напряжение от ТН, установленного на ошиновке НН или на вводе секции. Порог срабатывания по напряжению задается уставкой \textbf{«Uл КОН НН»}. Ввод в работу КОН НН осуществляется с помощью программного переключателя \textbf{«КОН НН»}. 

\item
При появлении неисправности в цепях напряжения стороны НН, КОН НН блокируется до исчезновения неисправности.  

\item
Оперативный ввод/вывод КОН НН производится с помощью БУ (блоков управления функциями) \textit{«Вывод КОН НН»}. 

\item
Орган, контролирующий отсутствие напряжения своей стороны (КОН СШ), контролирует минимальное линейное напряжение от ТН, установленного на системе шин своей стороны. Порог срабатывания по напряжению задается уставкой \textbf{«Uл КОН СШ»}. Ввод в работу КОН СШ осуществляется с помощью программного переключателя \textbf{«КОН СШ»}.

\item
При появлении неисправности в цепях напряжения своей стороны, КОН СШ блокируется до исчезновения неисправности.  

\item
Оперативный ввод/вывод КОН СШ производится с помощью БУ (блоков управления функциями) \textit{«Вывод КОН СШ»}. 


\end{enumerate}